\documentclass{report}

\usepackage[spanish]{babel}
\usepackage[utf8x]{inputenc}
\usepackage{amsmath}
\usepackage{graphicx}
\usepackage[colorinlistoftodos]{todonotes}
\usepackage{enumitem}
\usepackage{listings}
\usepackage{verbatim}
\usepackage{eurosym}
\usepackage[export]{adjustbox}
\usepackage{amssymb}
\usepackage{bussproofs}
\usepackage{amsmath}
\usepackage{tikz}
\usepackage{xcolor}
\usepackage{listings}
\usepackage{titletoc}
\usepackage{hyperref}

\hypersetup{
  colorlinks=true,
  linkcolor=black,
  urlcolor=blue,
  citecolor=black
}

\newcommand{\coverPage}[6]{%
%----------------------------------------------------------------------------------------
%	COVER START
%----------------------------------------------------------------------------------------
\begin{titlepage}

    \newcommand{\HRule}{\rule{\linewidth}{0.5mm}}
    \newcommand{\department}{#1}
    \newcommand{\course}{#2}
    \newcommand{\titleValue}{#3}
    \newcommand{\subtitleValue}{#4}
    \newcommand{\authorName}{#5}

    \center

    %----------------------------------------------------------------------------------------
    %	HEADER
    %----------------------------------------------------------------------------------------
    \includegraphics{images/logo_usa.png}
    \vspace{0.5cm}
    \textsc{\Large \department}\\[0.5cm]
    \textsc{\Large \course}\\[0.5cm]
    \vfill

    %----------------------------------------------------------------------------------------
    %	TITLE
    %----------------------------------------------------------------------------------------

    \HRule\\
    \Huge
    \textbf{\titleValue}\\[0.5cm]
    \Large
    \textbf{\subtitleValue}\\
    \HRule\\[0.5cm]

    %----------------------------------------------------------------------------------------
    %	AUTHOR AND DATE
    %----------------------------------------------------------------------------------------

    \vfill
    \Large
    \textit{\authorName}\\
    {\large \today}\\[2cm]

\end{titlepage}
%----------------------------------------------------------------------------------------
%	COVER END
%----------------------------------------------------------------------------------------
}

\begin{document}
    \coverPage{ Matemáticas }{ Geometría Euclidiana }{ Examen Final }{   }{ Alexander Mendoza }{\today}
    \tableofcontents

    \pagebreak
    \chapter{ Examen Final }

    \begin{enumerate}
        \setcounter{enumi}{1}
        \item Sea el $\triangle A B C, D$ un punto tal que $B-D-C$ y $E$ otro punto tal que $A-E-D$. Demostrar que $m \angle A E B>m \angle A C B$.

        \begin{center}
            \includegraphics*[height=5cm]{images/figura1.png}
        \end{center}
        \textit{\textbf{Demostración}}. Sea el $\triangle A B C, D$ un punto tal que $B-D-C$ y $E$ otro punto tal que $A-E-D$. Luego el $\angle AEB$ es externo al $\triangle BED$, en particular $m\angle AEB > m\angle ADB$, esto por el teorema del ángulo externo. De manera similar, el $\angle ADB$ es externo al $\triangle DAC$ y en particular el $m\angle ADB > m\angle ACB$. Por transitividad, $m\angle AEB > m\angle ACB$.


        \item El $\triangle P Q R$ es isósceles, la bisectriz de uno de los ángulos de la base, $\angle Q$, interseca al lado opuesto en un punto $S . T$ es un punto en la base $\overline{P Q}$ tal que $S T=P T$. La bisectriz del $\angle P S T$ interseca al $\overline{P Q}$ en un punto $V$. Demostrar que $\angle T S V \cong \angle R Q S$.

        \begin{center}
            \includegraphics*[height=5cm]{images/figura2.png}
        \end{center}
        \textit{\textbf{Demostración}}. Sea $\triangle P Q R$ un triángulo isosceles, la bisectriz de uno de los ángulos de la base, $\angle Q$, interseca al lado opuesto en un punto $S . T$ es un punto en la base $\overline{P Q}$ tal que $S T=P T$. La bisectriz del $\angle P S T$ interseca al $\overline{P Q}$ en un punto $V$. Luego como $T$ equidista de $S$ y $P$, el $\triangle SPT$ es isósceles con base $\overline{SP}$, esto por definición de triángulo isósceles. Luego $\angle TSP \cong \angle TPS$, como $S \in \overline{RP}$ y $T \in \overline{PQ}$, $\angle QPR = \angle TPS$. Luego como $\triangle QRP$ es isósceles, $\angle QPR \cong \angle PQR$, por transitividad $\angle TSP \cong \angle PQR$. De esta manera, $m\angle TSP = m\angle PQR$, como la mediatriz divide el ángulo en dos ángulos iguales tenemos que $\angle TSV \cong RQS$.

    \end{enumerate}
\end{document}
